\documentclass[11pt,a4paper]{article}

% Paquetes esenciales
\usepackage[utf8]{inputenc}
\usepackage[spanish]{babel}
\usepackage[margin=2.5cm]{geometry}
\usepackage{tabularx}
\usepackage{xcolor}
\usepackage{colortbl}
\usepackage{enumitem}
\usepackage{titlesec}
\usepackage{helvet}

% Configuración de fuente y formato
\renewcommand{\familydefault}{\sfdefault}
\definecolor{primary}{RGB}{44, 62, 80}
\definecolor{critical}{RGB}{192, 57, 43}
\definecolor{high}{RGB}{211, 84, 0}
\definecolor{graybg}{RGB}{242, 242, 242}

\titleformat{\section}{\Large\bfseries\color{primary}}{\thesection}{1em}{}[\titlerule]
\titleformat{\subsection}{\large\bfseries\color{primary}}{\thesubsection}{1em}{}

\begin{document}

% Encabezado
\begin{center}
    {\Huge \textbf{\color{primary}Reporte Ejecutivo de Riesgos Tecnológicos}}\\[0.5cm]
    {\large \textbf{Destinatario:} Directorio de la Clínica}\\[0.2cm]
    {\large \textbf{Fecha:} 5 de Septiembre de 2026}\\[0.5cm]
\end{center}

\vspace{0.5cm}

% 1. Resumen Ejecutivo
\section*{1. Resumen Ejecutivo}
El presente informe detalla el estado actual de la postura de ciberseguridad y continuidad operativa de la clínica. Tras una evaluación exhaustiva de los activos tecnológicos, se identificaron vulnerabilidades críticas que exponen a la institución a interrupciones severas en la atención de pacientes, pérdidas financieras directas y sanciones legales por incumplimiento de la Ley 25.326 de Protección de los Datos Personales. La estrategia prioritaria exige la adopción inmediata de medidas de mitigación tecnológicas (Backups inmutables, MFA) y procedimentales (Control de accesos) para neutralizar los vectores de ataque más probables, garantizando la integridad de las historias clínicas y la facturación en tiempo real.

\vspace{0.5cm}

% 2. Inventario de Activos
\section*{2. Inventario y Clasificación de Activos Críticos}
A continuación se detallan los activos tecnológicos fundamentales afectados por los riesgos identificados, clasificados según su nivel de criticidad.

\vspace{0.3cm}
\noindent
\begin{tabularx}{\textwidth}{|c|X|c|c|c|c|}
\hline
\rowcolor{primary} \textbf{\color{white}ID} & \textbf{\color{white}Descripción del Activo} & \textbf{\color{white}Tipo} & \textbf{\color{white}Responsable} & \textbf{\color{white}Clasificación} & \textbf{\color{white}Criticidad} \\ \hline
A01 & Base de datos de Historias Clínicas & Información & Jefatura Médica & Confidencial & Alta \\ \hline
A02 & Sistema de Facturación & Software & Ger. Financiera & Restringida & Alta \\ \hline
A03 & Personal del área de Facturación & Humano & Ger. Financiera & Interna & Media \\ \hline
A04 & Sistema de Gestión Hospitalaria (HIS) & Software & Gerencia TI & Confidencial & Alta \\ \hline
A05 & Enlace principal a Internet (ISP) & Red & Gerencia TI & Pública & Alta \\ \hline
A06 & Servidores físicos del Data Center & Hardware & Gerencia TI & Restringida & Alta \\ \hline
\end{tabularx}

\vspace{0.8cm}

% 3. Top 5 Riesgos por Nivel
\section*{3. Top 5 Riesgos por Nivel}
Los siguientes riesgos representan la mayor amenaza para la organización, calculados en base a su impacto (1-5) y probabilidad (1-5).

\vspace{0.3cm}
\noindent
\begin{tabularx}{\textwidth}{|c|X|c|c|c|}
\hline
\rowcolor{primary} \textbf{\color{white}\#} & \textbf{\color{white}Riesgo Identificado} & \textbf{\color{white}Impacto} & \textbf{\color{white}Probabilidad} & \textbf{\color{white}Puntaje} \\ \hline
1 & \textbf{Infección por Ransomware} en bases de datos. & 5 & 4 & \cellcolor{critical!30}\textbf{20 (Crítico)} \\ \hline
2 & \textbf{Fraude financiero por Phishing} a facturación. & 4 & 4 & \cellcolor{critical!30}\textbf{16 (Crítico)} \\ \hline
3 & \textbf{Exposición de datos} por exceso de privilegios. & 4 & 3 & \cellcolor{high!30}\textbf{12 (Alto)} \\ \hline
4 & \textbf{Interrupción de Obras Sociales} por corte de ISP. & 4 & 3 & \cellcolor{high!30}\textbf{12 (Alto)} \\ \hline
5 & \textbf{Apagón del Data Center} por falla eléctrica. & 5 & 2 & \cellcolor{high!30}\textbf{10 (Alto)} \\ \hline
\end{tabularx}

\vspace{0.8cm}

% 4. Estado de los Planes de Acción
\section*{4. Estado de los Planes de Acción (Riesgos Altos/Críticos)}
Se han diagramado tres planes de mitigación inmediatos para contener las amenazas más urgentes.

\vspace{0.3cm}
\begin{itemize}[leftmargin=*]
    \item \textbf{Plan 1: Arquitectura de Backups Inmutables y EDR (Ransomware)}
    \begin{itemize}
        \item \textbf{Estado:} Planificado.
        \item \textbf{Vencimiento:} 30/10/2026.
        \item \textbf{Presupuesto:} \$5,500 USD.
        \item \textbf{Responsable:} Gerencia de TI.
    \end{itemize}
    \vspace{0.2cm}
    
    \item \textbf{Plan 2: MFA Obligatorio y Concientización (Phishing)}
    \begin{itemize}
        \item \textbf{Estado:} En curso (Evaluación de proveedores).
        \item \textbf{Vencimiento:} 15/10/2026.
        \item \textbf{Presupuesto:} \$1,200 USD.
        \item \textbf{Responsable:} Gerencia Financiera.
    \end{itemize}
    \vspace{0.2cm}
    
    \item \textbf{Plan 3: Transición a RBAC en Recepción (Exposición de Datos)}
    \begin{itemize}
        \item \textbf{Estado:} Planificado.
        \item \textbf{Vencimiento:} 15/11/2026.
        \item \textbf{Presupuesto:} \$0 (Costos operativos internos).
        \item \textbf{Responsable:} Jefatura de Recepción.
    \end{itemize}
\end{itemize}

\vspace{0.5cm}

% 5. Recomendaciones Prioritarias
\section*{5. Recomendaciones Prioritarias}
Para asegurar la viabilidad operativa de la clínica en el corto plazo, se emiten las siguientes recomendaciones vinculantes al Directorio:

\begin{enumerate}
    \item \textbf{Aprobación Presupuestaria Inmediata:} Liberar los fondos estimados (\$6,700 USD) para ejecutar los Planes 1 y 2, cerrando las brechas críticas que actualmente permiten la ejecución de secuestros de datos y estafas financieras.
    \item \textbf{Adopción de Políticas Zero Trust:} Eliminar definitivamente el uso de cuentas genéricas compartidas en los sistemas de gestión hospitalaria (HIS), exigiendo responsabilidad individual en el manejo de datos de pacientes para cumplir con normativas legales.
    \item \textbf{Redundancia de Infraestructura Crítica:} Iniciar la contratación de un Proveedor de Servicios de Internet (ISP) secundario para garantizar la facturación y validación ininterrumpida con las Obras Sociales.
\end{enumerate}

\end{document}